\ifdefined\gkver\else\def\gkver{student}\fi
\documentclass[\gkver]{gaokaozhenti}
\begin{document}

\chapter{2020年海南}

\begin{enumerate}
\item
%% number: 1
%% paperName: 2020年海南高考第 1 题 3 分
%% typeId: 1
%% type: 单选题
%% chapter: 原子和原子核
%% point: 人工核反应
%% method: 无
%% score: 3
%% degree: 950
%% duplicateId: 0
%% body:
100 年前，卢瑟福猜想在原子核内除质子外还存在着另一种粒子 X，后来科学家用 $\alpha$ 粒子轰击铍核证实了这一猜想，该核反应方程为：\ce{^{4}_{2}He + ^{9}_{4}Be -> ^{12}_{6}C + ^{m}_{n}X}，则\xzanswer{A}
\fourchoices[answer=A]
{$m=1$，$n=0$，X 是中子}
{$m=1$，$n=0$，X 是电子}
{$m=0$，$n=1$，X 是中子}
{$m=0$，$n=1$，X 是电子}
%% answer: A
%% memo:
\memoanswer{
根据电荷数和质量数守恒，则有

$4+9=12+m$，$2+4=6+n$

解得 $m=1$，$n=0$，故 X 是中子。

故选 A。
}

\item
%% number: 2
%% paperName: 2020年海南高考第 2 题 3 分
%% typeId: 1
%% type: 单选题
%% chapter: 力物体的平衡
%% point: 三力平衡
%% method: 无
%% score: 3
%% degree: 850
%% duplicateId: 0
%% body:
如图，上网课时小明把手机放在斜面上，手机处于静止状态。则斜面对手机的\xzanswer{B}
\onepicture[w=3.5cm]{figs/02.png}
\fourchoices[answer=B]
{支持力竖直向上}
{支持力小于手机所受的重力}
{摩擦力沿斜面向下}
{摩擦力大于手机所受的重力沿斜面向下的分力}
%% answer: B
%% memo:
\memoanswer{
设手机的质量为 $m$，斜面倾角为 $\theta$。对手机进行受力分析，如图所示

\onepicture[w=4.0cm]{figs/02_an.png}

由图可知支持力方向垂直斜面向上，摩擦力方向沿斜面向上，根据平衡条件则有

$f=mg\sin\theta$，$F_{\mathrm{N}}=mg\cos\theta$

因 $\cos\theta<1$，故 $F_{\mathrm{N}}<mg$，且摩擦力等于手机所受的重力沿斜面向下的分力。

故选 B。
}

\item
%% number: 3
%% paperName: 2020年海南高考第 3 题 3 分
%% typeId: 1
%% type: 单选题
%% chapter: 交流电
%% point: 交流电
%% method: 无
%% score: 3
%% degree: 850
%% duplicateId: 0
%% body:
图甲、乙分别表示两种电流的波形，其中图乙所示电流按正弦规律变化，分别用 $I_{1}$ 和 $I_{2}$ 表示甲和乙两电流的有效值，则\xzanswer{D}
\twopicture[numstyle=chinese, wA=5.5cm, wB=5.5cm]{figs/03a.png}{figs/03b.png}
\fourchoices[answer=D]
{$I_{1}:I_{2}=2:1$}
{$I_{1}:I_{2}=1:2$}
{$I_{1}:I_{2}=1:\sqrt{2}$}
{$I_{1}:I_{2}=\sqrt{2}:1$}
%% answer: D
%% memo:
\memoanswer{
对图甲的交流电分析，可知一个周期内交流电的电流方向变化，而电流的大小不变，故图甲的电流有效值为 $I_{1}=I_{0}$；对图乙的交流电分析可知，其为正弦式交流电，故其有效值为 $I_{2}=\frac{I_{0}}{\sqrt{2}}$，故 $I_{1}:I_{2}=\sqrt{2}:1$。

故选 D。
}

\item
%% number: 4
%% paperName: 2020年海南高考第 4 题 3 分
%% typeId: 1
%% type: 单选题
%% chapter: 电路
%% point: 电功电功率
%% method: 无
%% score: 3
%% degree: 750
%% duplicateId: 0
%% body:
一车载加热器（额定电压为 $24\UV$）发热部分的电路如图所示，a、b、c 是三个接线端点，设 ab、ac、bc 间的功率分别为 $P_{ab}$、$P_{ac}$、$P_{bc}$，则\xzanswer{D}
\onepicture[w=3.5cm]{figs/04.png}
\fourchoices[answer=D]
{$P_{ab}>P_{bc}$}
{$P_{ab}=P_{ac}$}
{$P_{ac}=P_{bc}$}
{$P_{ab}<P_{ac}$}
%% answer: D
%% memo:
\memoanswer{
接 ab，则电路的总电阻为

$R_{ab}=\frac{9R(R+9R)}{R+9R+9R}=\frac{90R}{19}$

接 ac，则电路的总电阻为

$R_{ac}=\frac{R(9R+9R)}{R+9R+9R}=\frac{18R}{19}$

接 bc，则电路的总电阻为

$R_{bc}=\frac{9R(R+9R)}{R+9R+9R}=\frac{90R}{19}$

由题知，不管接哪两个点，电压不变，为 $U=24\UV$，根据 $P=\frac{U^{2}}{R}$ 可知 $P_{ab}=P_{bc}<P_{ac}$。

故选 D。
}

\item
%% number: 5
%% paperName: 2020年海南高考第 5 题 3 分
%% typeId: 1
%% type: 单选题
%% chapter: 光学
%% point: 光的电磁说
%% method: 无
%% score: 3
%% degree: 750
%% duplicateId: 0
%% body:
下列说法正确的是\xzanswer{A}
\fourchoices[answer=A]
{单色光在介质中传播时，介质的折射率越大，光的传播速度越小}
{观察者靠近声波波源的过程中，接收到的声波频率小于波源频率}
{同一个双缝干涉实验中，蓝光产生的干涉条纹间距比红光的大}
{两束频率不同的光，可以产生干涉现象}
%% answer: A
%% memo:
\memoanswer{
A．根据 $v=\frac{c}{n}$ 可知，单色光在介质中传播时，介质的折射率越大，光的传播速度越小，故 A 正确；

B．根据多普勒效应，若声波波源向观察者靠近，则观察者接收到的声波频率大于波源频率，故 B 错误；

C．根据 $\Delta x=\frac{L}{d}\lambda$，同一个双缝干涉实验中，蓝光的波长小于红光的波长，故蓝光产生的干涉条纹间距比红光的小，故 C 错误；

D．根据光的干涉的条件可知，两束频率不同的光不能产生干涉现象，故 D 错误。

故选 A。
}

\item
%% number: 6
%% paperName: 2020年海南高考第 6 题 3 分
%% typeId: 1
%% type: 单选题
%% chapter: 磁场
%% point: 左手定则
%% method: 无
%% score: 3
%% degree: 850
%% duplicateId: 0
%% body:
如图，在一个蹄形电磁铁的两个磁极的正中间放置一根长直导线，当导线中通有垂直于纸面向里的电流 $I$ 时，导线所受安培力的方向为\xzanswer{B}
\onepicture[w=3.2cm]{figs/06.png}
\fourchoices[answer=B]
{向上}
{向下}
{向左}
{向右}
%% answer: B
%% memo:
\memoanswer{
根据安培定则，可知蹄形电磁铁的磁极分布情况，如图所示

\onepicture[w=3.2cm]{figs/06_an.png}

故导线所处位置的磁感应线的切线方向为水平向右，根据左手定则，可以判断导线所受安培力的方向为向下。

故选 B。
}

\item
%% number: 7
%% paperName: 2020年海南高考第 7 题 3 分
%% typeId: 1
%% type: 单选题
%% chapter: 曲线运动
%% point: 人造地球卫星
%% method: 无
%% score: 3
%% degree: 550
%% duplicateId: 0
%% body:
2020 年 5 月 5 日，长征五号 B 运载火箭在中国文昌航天发射场成功首飞，将新一代载人飞船试验船送入太空，若试验船绕地球做匀速圆周运动，周期为 $T$，离地高度为 $h$，已知地球半径为 $R$，万有引力常量为 $G$，则\xzanswer{B}
\fourchoices[answer=B]
{试验船的运行速度为 $\frac{2\pi R}{T}$}
{地球的第一宇宙速度为 $\frac{2\pi}{T}\sqrt{\frac{(R+h)^{3}}{R}}$}
{地球的质量为 $\frac{2\pi(R+h)^{3}}{GT^{2}}$}
{地球表面的重力加速度为 $\frac{4\pi^{2}(R+h)^{2}}{RT^{2}}$}
%% answer: B
%% memo:
\memoanswer{
A．试验船的运行速度为 $\frac{2\pi(R+h)}{T}$，故 A 错误；

B．近地轨道卫星的速度等于第一宇宙速度，根据万有引力提供向心力有 $G\frac{Mm}{R^{2}}=m\frac{v^{2}}{R}$；根据试验船受到的万有引力提供向心力有 $G\frac{Mm_{\text{船}}}{(R+h)^{2}}=m_{\text{船}}\left(\frac{2\pi}{T}\right)^{2}(R+h)$，联立两式解得第一宇宙速度 $v=\frac{2\pi}{T}\sqrt{\frac{(R+h)^{3}}{R}}$，故 B 正确；

C．根据试验船受到的万有引力提供向心力有 $G\frac{Mm_{\text{船}}}{(R+h)^{2}}=m_{\text{船}}\left(\frac{2\pi}{T}\right)^{2}(R+h)$，解得 $M=\frac{4\pi^{2}(R+h)^{3}}{GT^{2}}$，故 C 错误；

D．地球重力加速度等于近地轨道卫星向心加速度，根据万有引力提供向心力有 $G\frac{Mm}{R^{2}}=m\frac{v^{2}}{R}=mg$；根据试验船受到的万有引力提供向心力有 $G\frac{Mm_{\text{船}}}{(R+h)^{2}}=m_{\text{船}}\left(\frac{2\pi}{T}\right)^{2}(R+h)$，联立两式解得重力加速度 $g=\frac{4\pi^{2}(R+h)^{3}}{R^{2}T^{2}}$，故 D 错误。

故选 B。
}

\item
%% number: 8
%% paperName: 2020年海南高考第 8 题 3 分
%% typeId: 1
%% type: 单选题
%% chapter: 运动和力
%% point: 动量定理
%% method: 无
%% score: 3
%% degree: 650
%% duplicateId: 0
%% body:
太空探测器常装配离子发动机，其基本原理是将被电离的原子从发动机尾部高速喷出，从而为探测器提供推力，若某探测器质量为 $490\Ukg$，离子以 $30\Ukms$ 的速率（远大于探测器的飞行速率）向后喷出，流量为 $3.0\times10^{-3}\Ugs$，则探测器获得的平均推力大小为\xzanswer{C}
\fourchoices[answer=C]
{$1.47\UN$}
{$0.147\UN$}
{$0.09\UN$}
{$0.009\UN$}
%% answer: C
%% memo:
\memoanswer{
对离子，根据动量定理有 $F\cdot\Delta t=\Delta mv$，而 $\Delta m=3.0\times10^{-3}\times10^{-3}\Delta t\Ukg$，解得 $F=0.09\UN$，故探测器获得的平均推力大小为 $0.09\UN$。

故选 C。
}

\item
%% number: 9
%% paperName: 2020年海南高考第 9 题 4 分
%% typeId: 2
%% type: 多选题
%% chapter: 机械波
%% point: 波的图象
%% method: 无
%% score: 4
%% degree: 750
%% duplicateId: 0
%% body:
一列简谐横波沿 $x$ 轴正方向传播，波的周期为 $0.2\Us$，某时刻的波形如图所示。则\xzanswer{AC}
\onepicture[w=6.0cm]{figs/09.png}
\fourchoices[answer=AC]
{该波的波长为 $8\Um$}
{该波的波速为 $50\Ums$}
{该时刻质点 P 向 $y$ 轴负方向运动}
{该时刻质点 Q 向 $y$ 轴负方向运动}
%% answer: AC
%% memo:
\memoanswer{
A．由波形图可知，波长为 $8\Um$，故 A 正确；

B．根据公式 $v=\frac{\lambda}{T}$，代入数据解得 $v=40\Ums$，故 B 错误；

C、D．由题知，沿 $x$ 轴正方向传播，根据“上下坡法”，可知该时刻质点 P 向 $y$ 轴负方向运动，该时刻质点 Q 向 $y$ 轴正方向运动，故 C 正确，D 错误。

故选 AC。
}

\item
%% number: 10
%% paperName: 2020年海南高考第 10 题 4 分
%% typeId: 2
%% type: 多选题
%% chapter: 电场
%% point: 电场线
%% method: 无
%% score: 4
%% degree: 650
%% duplicateId: 0
%% body:
空间存在如图所示的静电场，a、b、c、d 为电场中的四个点，则\xzanswer{AD}
\onepicture[w=3.5cm]{figs/10.png}
\fourchoices[answer=AD]
{a 点的场强比 b 点的大}
{d 点的电势比 c 点的低}
{质子在 d 点的电势能比在 c 点的小}
{将电子从 a 点移动到 b 点，电场力做正功}
%% answer: AD
%% memo:
\memoanswer{
A．根据电场线的疏密程度表示电场强度的大小，可知 a 点的电场线比 b 点的电场线更密，故 a 点的场强比 b 点的场强大，故 A 正确；

B．根据沿着电场线方向电势不断降低，可知 d 点的电势比 c 点的电势高，故 B 错误；

C．根据正电荷在电势越高的点电势能越大，可知质子在 d 点的电势能比在 c 点的电势能大，故 C 错误；

D．由图可知，a 点的电势低于 b 点的电势，而负电荷在电势越低的点电势能越大，故电子在 a 点的电势能高于在 b 点的电势能，所以将电子从 a 点移动到 b 点，电势能减小，故电场力做正功，故 D 正确。

故选 AD。
}

\item
%% number: 11
%% paperName: 2020年海南高考第 11 题 4 分
%% typeId: 2
%% type: 多选题
%% chapter: 曲线运动
%% point: 平抛运动
%% method: 无
%% score: 4
%% degree: 650
%% duplicateId: 0
%% body:
小朋友玩水枪游戏时，若水从枪口沿水平方向射出的速度大小为 $10\Ums$，水射出后落到水平地面上。已知枪口离地高度为 $1.25\Um$，$g=10\Umsq$，忽略空气阻力，则射出的水\xzanswer{BD}
\fourchoices[answer=BD]
{在空中的运动时间为 $0.25\Us$}
{水平射程为 $5\Um$}
{落地时的速度大小为 $15\Ums$}
{落地时竖直方向的速度大小为 $5\Ums$}
%% answer: BD
%% memo:
\memoanswer{
A．根据 $h=\frac{1}{2}gt^{2}$ 得，运动时间 $t=\sqrt{\frac{2h}{g}}=\sqrt{\frac{2\times1.25}{10}}\Us=0.5\Us$，故 A 错误；

B．水平射程为 $x=v_{0}t=10\Ums\times0.5\Us=5\Um$，故 B 正确；

C、D．竖直方向分速度为 $v_{y}=gt=10\Umsq\times0.5\Us=5\Ums$，水平分速度为 $v_{x}=v_{0}=10\Ums$，落地速度为 $v=\sqrt{v_{y}^{2}+v_{x}^{2}}=5\sqrt{5}\Ums$，故 C 错误，D 正确。

故选 BD。
}

\item
%% number: 12
%% paperName: 2020年海南高考第 12 题 4 分
%% typeId: 2
%% type: 多选题
%% chapter: 运动和力
%% point: 连接体
%% method: 无
%% score: 4
%% degree: 450
%% duplicateId: 0
%% body:
如图，在倾角为 $\theta$ 的光滑斜面上，有两个物块 P 和 Q，质量分别为 $m_{1}$ 和 $m_{2}$，用与斜面平行的轻质弹簧相连接，在沿斜面向上的恒力 $F$ 作用下，两物块一起向上做匀加速直线运动，则\xzanswer{BC}
\onepicture[w=4.0cm]{figs/12.png}
\fourchoices[answer=BC]
{两物块一起运动的加速度大小为 $a=\frac{F}{m_{1}+m_{2}}$}
{弹簧的弹力大小为 $T=\frac{m_{2}}{m_{1}+m_{2}}F$}
{若只增大 $m_{2}$，两物块一起向上匀加速运动时，它们的间距变大}
{若只增大 $\theta$，两物块一起向上匀加速运动时，它们的间距变大}
%% answer: BC
%% memo:
\memoanswer{
A．对整体受力分析，根据牛顿第二定律有 $F-(m_{1}+m_{2})g\sin\theta=(m_{1}+m_{2})a$，解得 $a=\frac{F}{m_{1}+m_{2}}-g\sin\theta$，故 A 错误；

B．对 $m_{2}$ 受力分析，根据牛顿第二定律有 $F_{\text{弹}}-m_{2}g\sin\theta=m_{2}a$，解得 $F_{\text{弹}}=\frac{m_{2}F}{m_{1}+m_{2}}$，故 B 正确；

C．$F_{\text{弹}}=\frac{m_{2}F}{m_{1}+m_{2}}=\frac{F}{\frac{m_{1}}{m_{2}}+1}$，可知若只增大 $m_{2}$，两物块一起向上匀加速运动时，弹力变大，根据胡克定律，可知伸长量变大，故它们的间距变大，故 C 正确；

D．根据 $F_{\text{弹}}=\frac{m_{2}F}{m_{1}+m_{2}}$，可知只增大 $\theta$，两物块一起向上匀加速运动时，弹力不变，根据胡克定律，可知伸长量不变，故它们的间距不变，故 D 错误。

故选 BC。
}

\item
%% number: 13
%% paperName: 2020年海南高考第 13 题 4 分
%% typeId: 2
%% type: 多选题
%% chapter: 电磁感应
%% point: 电磁感应能量分析
%% method: 微元
%% score: 4
%% degree: 350
%% duplicateId: 0
%% body:
如图，足够长的间距 $d=1\Um$ 的平行光滑金属导轨 MN、PQ 固定在水平面内，导轨间存在一个宽度 $L=1\Um$ 的匀强磁场区域，磁感应强度大小为 $B=0.5\UT$，方向如图所示。一根质量 $m_{a}=0.1\Ukg$，阻值 $R=0.5\UO$ 的金属棒 a 以初速度 $v_{0}=4\Ums$ 从左端开始沿导轨滑动，穿过磁场区域后，与另一根质量 $m_{b}=0.2\Ukg$，阻值 $R=0.5\UO$ 的原来静置在导轨上的金属棒 b 发生弹性碰撞，两金属棒始终与导轨垂直且接触良好，导轨电阻不计，则\xzanswer{BD}
\onepicture[w=6.0cm]{figs/13.png}
\fourchoices[answer=BD]
{金属棒 a 第一次穿过磁场时做匀减速直线运动}
{金属棒 a 第一次穿过磁场时回路中有逆时针方向的感应电流}
{金属棒 a 第一次穿过磁场区域的过程中，金属棒 b 上产生的焦耳热为 $0.25\UJ$}
{金属棒 a 最终停在距磁场左边界 $0.8\Um$ 处}
%% answer: BD
%% memo:
\memoanswer{
A．金属棒 a 第一次穿过磁场时受到安培力的作用，做减速运动，由于速度减小，感应电流减小，安培力减小，加速度减小，故金属棒 a 做加速度减小的减速直线运动，故 A 错误；

B．根据右手定则可知，金属棒 a 第一次穿过磁场时回路中有逆时针方向的感应电流，故 B 正确；

C．电路中产生的平均电动势为 $\overline{E}=\frac{\Delta\Phi}{\Delta t}=\frac{BLd}{\Delta t}$，平均电流为 $\overline{I}=\frac{\overline{E}}{2R}$，金属棒 a 受到的安培力为 $\overline{F}=B\overline{I}d$。规定向右为正方向，对金属棒 a，根据动量定理得 $-B\overline{I}d\cdot\Delta t=m_{a}v_{a}-m_{a}v_{0}$，解得金属棒第一次离开磁场时速度 $v_{a}=1.5\Ums$。金属棒 a 第一次穿过磁场区域的过程中，电路中产生的总热量等于金属棒 a 机械能的减少量，即 $Q=\frac{1}{2}m_{a}v_{0}^{2}-\frac{1}{2}m_{a}v_{a}^{2}$，联立并代入数据得 $Q=0.6875\UJ$。由于两棒电阻相同，两棒产生的焦耳热相同，则金属棒 b 上产生的焦耳热 $Q_{b}=\frac{Q}{2}=0.34375\UJ$，故 C 错误；

D．规定向右为正方向，两金属棒碰撞过程根据动量守恒和机械能守恒得 $m_{a}v_{a}=m_{a}v_{a}'+m_{b}v_{b}$，$\frac{1}{2}m_{a}v_{a}^{2}=\frac{1}{2}m_{a}v_{a}'^{2}+\frac{1}{2}m_{b}v_{b}^{2}$，联立并代入数据解得金属棒 a 反弹的速度为 $v_{a}'=-0.5\Ums$。设金属棒 a 最终停在距磁场左边界 $x$ 处，则从反弹进入磁场到停下来的过程，电路中产生的平均电动势为 $\overline{E}'=\frac{\Delta\Phi'}{\Delta t'}=\frac{B(L-x)d}{\Delta t'}$，平均电流为 $\overline{I}'=\frac{\overline{E}'}{2R}$，金属棒 a 受到的安培力为 $\overline{F}'=B\overline{I}'d$，规定向右为正方向，对金属棒 a，根据动量定理得 $-B\overline{I}'d\cdot\Delta t=0-m_{a}v_{a}'$，联立并代入数据解得 $x=0.8\Um$，故 D 正确。

故选 BD。
}

\item
%% number: 14
%% paperName: 2020年海南高考第 14 题 6 分
%% typeId: 4
%% type: 实验题
%% chapter: 直线运动
%% point: 自由落体运动
%% method: 无
%% score: 6
%% degree: 650
%% duplicateId: 0
%% body:
（1）滑板运动场地有一种常见的圆弧形轨道，其截面如图所示，某同学用一辆滑板车和手机估测轨道半径 $R$（滑板车的长度远小于轨道半径）。主要实验过程如下：
\begin{enumerate}
	\item
	用手机查得当地的重力加速度 $g$；
	\item
	找出轨道的最低点 O，把滑板车从 O 点移开一小段距离至 P 点，由静止释放，用手机测出它完成 $n$ 次全振动的时间 $t$，算出滑板车做往复运动的周期 $T=$ \tkanswer{$\frac{t}{n}$}；
	\item
	将滑板车的运动视为简谐运动，则可将以上测量结果代入公式 $R=$ \tkanswer{$\frac{gt^{2}}{4n^{2}\pi^{2}}$}（用 $T$、$g$ 表示）计算出轨道半径。
\end{enumerate}

\onepicture[align=right, w=7.0cm]{figs/14a.png}

（2）某同学用如图（a）所示的装置测量重力加速度。

实验器材：有机玻璃条（白色是透光部分，黑色是宽度均为 $d=1.00\Ucm$ 的挡光片），铁架台，数字计时器（含光电门），刻度尺。

主要实验过程如下：
\begin{enumerate}
	\item
	将光电门安装在铁架台上，下方放置承接玻璃条下落的缓冲物；
	\item
	用刻度尺测量两挡光片间的距离，刻度尺的示数如图（b）所示，读出两挡光片间的距离 $L=$ \tkanswer{15.40}~$\Ucm$；
	\item
	手提玻璃条上端使它静止在 \tkanswer{竖直} 方向上，让光电门的光束从玻璃条下端的透光部分通过；
	\item
	让玻璃条自由下落，测得两次挡光的时间分别为 $t_{1}=10.003$ ms 和 $t_{2}=5.000$ ms；
	\item
	根据以上测量的数据计算出重力加速度 $g=$ \tkanswer{9.74}~$\Umsq$（结果保留三位有效数字）。
\end{enumerate}

\twopicture[align=right, showlabel=false, wA=2.6cm, wB=7.0cm]{figs/14b.png}{figs/14c.png}
%% answer:
\jdanswer{
\begin{enumerate}
	\item
	$\frac{t}{n}$，$\frac{gt^{2}}{4n^{2}\pi^{2}}$
	\item
	15.40，竖直，9.74
\end{enumerate}
}
%% memo:
\memoanswer{
（1）滑板车做往复运动的周期为 $T=\frac{t}{n}$；根据单摆的周期公式 $T=2\pi\sqrt{\frac{R}{g}}$，得 $R=\frac{gt^{2}}{4\pi^{2}}=\frac{gt^{2}}{4n^{2}\pi^{2}}$。

（2）两挡光片间的距离 $L=15.40\Ucm-0\Ucm=15.40\Ucm$；手提玻璃条上端使它静止在竖直方向上，让光电门的光束从玻璃条下端的透光部分通过。玻璃条下部挡光条通过光电门时玻璃条的速度为 $v_{1}=\frac{d}{t_{1}}=\frac{1.00\times10^{-2}}{10.003\times10^{-3}}\Ums=1\Ums$；玻璃条上部挡光条通过光电门时玻璃条的速度为 $v_{2}=\frac{d}{t_{2}}=\frac{1.00\times10^{-2}}{5.000\times10^{-3}}\Ums=2\Ums$；根据速度位移公式有 $v_{2}^{2}-v_{1}^{2}=2gL$，代入数据解得加速度 $g=\frac{v_{2}^{2}-v_{1}^{2}}{2L}=9.74\Umsq$。
}

\item
%% number: 15
%% paperName: 2020年海南高考第 15 题 10 分
%% typeId: 4
%% type: 实验题
%% chapter: 电路
%% point: 伏安法测电阻
%% method: 无
%% score: 10
%% degree: 450
%% duplicateId: 0
%% body:
在测量定值电阻阻值的实验中，提供的实验器材如下：电压表 V$_{1}$（量程 $3\UV$，内阻 $r_{1}=3.0\UkO$），电压表 V$_{2}$（量程 $5\UV$，内阻 $r_{2}=5.0\UkO$），滑动变阻器 $R$（额定电流 $1.5\UA$，最大阻值 $100\UO$），待测定值电阻 $R_{x}$，电源 $E$（电动势 $6.0\UV$，内阻不计），单刀开关 S，导线若干：

回答下列问题：
\begin{enumerate}
	\item
	实验中滑动变阻器应采用 \tkanswer{分压} 接法（填“限流”或“分压”）；
	\item
	将虚线框中的电路原理图补充完整；
	\item
	根据下表中的实验数据（$U_{1}$、$U_{2}$ 分别为电压表 V$_{1}$、V$_{2}$ 的示数），在图（a）给出的坐标纸上补齐数据点，并绘制 $U_{2}-U_{1}$ 图像：

	\begin{center}
	\begin{tabular}{|c|c|c|c|c|c|}
	\hline
	测量次数 & 1 & 2 & 3 & 4 & 5 \\
	\hline
	$U_{1}/\UV$ & 1.00 & 1.50 & 2.00 & 2.50 & 3.00 \\
	\hline
	$U_{2}/\UV$ & 1.61 & 2.41 & 3.21 & 4.02 & 4.82 \\
	\hline
	\end{tabular}
	\end{center}
	\item
	由 $U_{2}-U_{1}$ 图像得到待测定值电阻的阻值 $R_{x}=$ \tkanswer{$1.83\times10^{3}$} $\UO$（结果保留三位有效数字）；
	\item
	完成上述实验后，若要继续采用该实验原理测定另一个定值电阻 $R_{y}$（阻值约为 $700\UO$）的阻值，在不额外增加器材的前提下，要求实验精度尽可能高，请在图（b）的虚线框内画出你改进的电路图。
\end{enumerate}

\twopicture[align=right, showlabel=false, wA=5.0cm, wB=5.0cm]{figs/15a.png}{figs/15b.png}
%% answer:
\jdanswer{
\begin{enumerate}
	\item
	分压
	\item
	如图甲所示
	\item
	如图乙所示
	\item
	$1.83\times10^{3}$
	\item
	如图丙所示
\end{enumerate}
}
%% memo:
\memoanswer{
（1）电压表 V$_{1}$ 和 $R_{x}$ 串联后再和 V$_{2}$ 并联，并联的总电阻远大于滑动变阻器电阻，为了调节变阻器时电表示数变化明显，且能获得更多组数据，选择分压接法。

（2）完整的电路图，如图甲所示。

\onepicture[num=甲, numstyle=chinese, w=4.5cm]{figs/15_an1.png}

（3）根据下表中的实验数据，绘制的 $U_{2}-U_{1}$ 图像，如图乙所示。

\onepicture[num=乙, numstyle=chinese, w=5.5cm]{figs/15_an2.png}

（4）根据实验电路图，则有 $R_{x}=\frac{U_{2}-U_{1}}{\frac{U_{1}}{r_{1}}}$，变形得 $U_{2}=\frac{R_{x}+r_{1}}{r_{1}}U_{1}$，则图线的斜率为 $k=\frac{R_{x}+r_{1}}{r_{1}}$。根据 $U_{2}-U_{1}$ 图像可得斜率 $k=\frac{4.82-1.61}{3.00-1.00}=1.61$，则有 $1.61=\frac{R_{x}+r_{1}}{r_{1}}$，代入 $r_{1}=3.0\UkO$，解得 $R_{x}=1.83\times10^{3}\UO$。

（5）因待测电阻 $R_{y}$（阻值约为 $700\UO$）的阻值较小，若仍与电压表 V$_{1}$ 串联，则所分得的电压过小，不利于测量，故待测电阻 $R_{x}$ 与 $R_{y}$ 串联后与其中一个电压表并联；由于电源电动势只有 $6\UV$，因两电阻之和与电压表 V$_{1}$ 的阻值相当，故让两电阻与电压表 V$_{1}$ 并联，再与电压表 V$_{2}$ 串联，故改进后的电路图如图丙所示。

\onepicture[num=丙, numstyle=chinese, w=4.5cm]{figs/15_an3.png}
}

\item
%% number: 16
%% paperName: 2020年海南高考第 16 题 8 分
%% typeId: 6
%% type: 计算题
%% chapter: 气体性质
%% point: 玻意耳定律
%% method: 无
%% score: 8
%% degree: 650
%% duplicateId: 0
%% body:
如图，圆柱形导热气缸长 $L_{0}=60\Ucm$，缸内用活塞（质量和厚度均不计）密闭了一定质量的理想气体，缸底装有一个触发器 D，当缸内压强达到 $p=1.5\times10^{5}\UPa$ 时，D 被触发，不计活塞与缸壁的摩擦。初始时，活塞位于缸口处，环境温度 $t_{0}=27$ ℃，压强 $p_{0}=1.0\times10^{5}\UPa$。
\begin{enumerate}
	\item
	若环境温度不变，缓慢向下推活塞，求 D 刚好被触发时，到缸底的距离；
	\item
	若活塞固定在缸口位置，缓慢升高环境温度，求 D 刚好被触发时的环境温度。
\end{enumerate}
\onepicture[align=right, w=3.2cm]{figs/16.png}
%% answer:
\jdanswer{
\begin{enumerate}
	\item
	$0.4\Um$
	\item
	$450\UK$
\end{enumerate}
}
%% memo:
\memoanswer{
（1）设气缸横截面积为 $S$；D 刚好被触发时，到缸底的距离为 $L$，根据玻意耳定律得 $p_{0}SL_{0}=pSL$，代入数据解得 $L=\frac{p_{0}}{p}L_{0}=\frac{1\times10^{5}}{1.5\times10^{5}}\times0.6\Um=0.4\Um$。

（2）此过程为等容变化，根据查理定律得 $\frac{p_{0}}{T_{0}}=\frac{p}{T}$，代入数据解得 $T=\frac{p}{p_{0}}T_{0}=\frac{1.5\times10^{5}}{1\times10^{5}}\times(27+273)\UK=450\UK$。
}

\item
%% number: 17
%% paperName: 2020年海南高考第 17 题 12 分
%% typeId: 6
%% type: 计算题
%% chapter: 运动和力
%% point: 动量和能量
%% method: 无
%% score: 12
%% degree: 550
%% duplicateId: 0
%% body:
如图，光滑的四分之一圆弧轨道 PQ 竖直放置，底端与一水平传送带相切，一质量 $m_{a}=1\Ukg$ 的小物块 a 从圆弧轨道最高点 P 由静止释放，到最低点 Q 时与另一质量 $m_{b}=3\Ukg$ 的小物块 b 发生弹性正碰（碰撞时间极短）。已知圆弧轨道半径 $R=0.8\Um$，传送带的长度 $L=1.25\Um$，传送带以速度 $v=1\Ums$ 顺时针匀速转动，小物块与传送带间的动摩擦因数 $\mu=0.2$，$g=10\Umsq$。求
\begin{enumerate}
	\item
	碰撞前瞬间小物块 a 对圆弧轨道的压力大小；
	\item
	碰后小物块 a 能上升的最大高度；
	\item
	小物块 b 从传送带的左端运动到右端所需要的时间。
\end{enumerate}
\onepicture[align=right, w=7.0cm]{figs/17.png}
%% answer:
\jdanswer{
\begin{enumerate}
	\item
	$30\UN$
	\item
	$0.2\Um$
	\item
	$1\Us$
\end{enumerate}
}
%% memo:
\memoanswer{
（1）设小物块 a 下到圆弧最低点未与小物块 b 相碰时的速度为 $v_{a}$，根据机械能守恒定律有 $m_{a}gR=\frac{1}{2}m_{a}v_{a}^{2}$，代入数据解得 $v_{a}=4\Ums$。小物块 a 在最低点，根据牛顿第二定律有 $F_{\mathrm{N}}-m_{a}g=m_{a}\frac{v_{a}^{2}}{R}$，代入数据解得 $F_{\mathrm{N}}=30\UN$。根据牛顿第三定律，可知小物块 a 对圆弧轨道的压力大小为 $30\UN$。

（2）小物块 a 与小物块 b 发生弹性碰撞，根据动量守恒有 $m_{a}v_{a}=m_{a}v_{a}'+m_{b}v_{b}$，根据能量守恒有 $\frac{1}{2}m_{a}v_{a}^{2}=\frac{1}{2}m_{a}v_{a}'^{2}+\frac{1}{2}m_{b}v_{b}^{2}$，联立解得 $v_{a}'=-2\Ums$，$v_{b}=2\Ums$。小物块 a 反弹，根据机械能守恒有 $m_{a}gh=\frac{1}{2}m_{a}v_{a}'^{2}$，解得 $h=0.2\Um$。

（3）小物块 b 滑上传送带，因 $v_{b}=2\Ums>v=1\Ums$，故小物块 b 先做匀减速运动，根据牛顿第二定律有 $\mu m_{b}g=m_{b}a$，解得 $a=2\Umsq$。则小物块 b 由 $2\Ums$ 减至 $1\Ums$，所走过的位移为 $x_{1}=\frac{v_{b}^{2}-v^{2}}{2a}$，代入数据解得 $x_{1}=0.75\Um$，运动的时间为 $t_{1}=\frac{v_{b}-v}{a}$，代入数据解得 $t_{1}=0.5\Us$。因 $x_{1}=0.75\Um<L=1.25\Um$，故小物块 b 之后将做匀速运动至右端，则匀速运动的时间为 $t_{2}=\frac{L-x_{1}}{v}=\frac{1.25-0.75}{1}\Us=0.5\Us$。故小物块 b 从传送带的左端运动到右端所需要的时间 $t=t_{1}+t_{2}=1\Us$。
}

\item
%% number: 18
%% paperName: 2020年海南高考第 18 题 16 分
%% typeId: 6
%% type: 计算题
%% chapter: 磁场
%% point: 带电粒子在磁场中的运动
%% method: 无
%% score: 16
%% degree: 350
%% duplicateId: 0
%% body:
如图，虚线 MN 左侧有一个正三角形 ABC，C 点在 MN 上，AB 与 MN 平行，该三角形区域内存在垂直于纸面向外的匀强磁场；MN 右侧的整个区域存在垂直于纸面向里的匀强磁场，一个带正电的离子（重力不计）以初速度 $v_{0}$ 从 AB 的中点 O 沿 OC 方向射入三角形区域，偏转 $60^\circ$ 后从 MN 上的 P 点（图中未画出）进入 MN 右侧区域，偏转后恰能回到 O 点。已知离子的质量为 $m$，电荷量为 $q$，正三角形的边长为 $d$。
\begin{enumerate}
	\item
	求三角形区域内磁场的磁感应强度；
	\item
	求离子从 O 点射入到返回 O 点所需要的时间；
	\item
	若原三角形区域存在的是一磁感应强度大小与原来相等的恒定磁场，将 MN 右侧磁场变为一个与 MN 相切于 P 点的圆形匀强磁场，让离子从 P 点射入圆形磁场，速度大小仍为 $v_{0}$，方向垂直于 BC，始终在纸面内运动，到达 O 点时的速度方向与 OC 成 $120^\circ$ 角，求圆形磁场的磁感应强度。
\end{enumerate}
\onepicture[align=right, w=3.2cm]{figs/18.png}
%% answer:
\jdanswer{
\begin{enumerate}
	\item
	$B=\frac{2mv_{0}}{qd}$
	\item
	$t=\frac{(11\pi+3\sqrt{3})d}{3v_{0}}$
	\item
	若三角形 ABC 区域磁场方向向里，$B_{2}=\frac{6mv_{0}}{5qd}$；若三角形 ABC 区域磁场方向向外，$B_{3}=\frac{2mv_{0}}{qd}$
\end{enumerate}
}
%% memo:
\memoanswer{
（1）画出粒子运动轨迹如图。

\onepicture[w=6.5cm]{figs/18_an1.png}

粒子在三角形 ABC 中运动时，有 $qBv_{0}=m\frac{v_{0}^{2}}{r}$，$T=\frac{2\pi r}{v_{0}}$，又粒子出三角形磁场时偏转 $60^\circ$，由几何关系可知 $r=\frac{d}{2}$，联立解得 $B=\frac{2mv_{0}}{qd}$，$t_{1}=\frac{T}{6}=\frac{\pi d}{6v_{0}}$。

（2）粒子从 D 运动到 P，由几何关系可知 $CP=d$，$DP=CP\sin60^\circ$，运动时间 $t_{2}=\frac{DP}{v_{0}}=\frac{\sqrt{3}d}{2v_{0}}$。粒子在 MN 右侧运动的半径为 $r'=2d$，则有 $qB'v_{0}=m\frac{v_{0}^{2}}{r'}$，$T'=\frac{2\pi r'}{v_{0}}$，运动时间 $t_{3}=\frac{5}{6}T'=\frac{10\pi d}{3v_{0}}$。故粒子从 O 点射入到返回 O 点所需要的时间 $t=2(t_{1}+t_{2})+t_{3}=\frac{11\pi d}{3v_{0}}+\frac{\sqrt{3}d}{v_{0}}$。

（3）若三角形 ABC 区域磁场方向向里，则粒子运动轨迹如图中①所示，有 $R+R\cos60^\circ=d+\frac{d}{2}\cos60^\circ$，解得 $R=\frac{5}{6}d$，此时根据 $qB_{2}v_{0}=m\frac{v_{0}^{2}}{R}$ 有 $B_{2}=\frac{6mv_{0}}{5qd}$；若三角形 ABC 区域磁场方向向外，则粒子运动轨迹如图中②所示，有 $R'+R'\cos60^\circ=d-\frac{d}{2}\cos60^\circ$，解得 $R'=\frac{1}{2}d$，此时根据 $qB_{3}v_{0}=m\frac{v_{0}^{2}}{R'}$ 有 $B_{3}=\frac{2mv_{0}}{qd}$。

\onepicture[w=7.0cm]{figs/18_an2.png}
}

\end{enumerate}

\end{document}
